\documentclass{article}
%packages
\usepackage{fullpage}
\usepackage{xspace}
\usepackage{graphicx,color}
\usepackage{hyperref}
\usepackage{graphicx}
\usepackage[dvipsnames]{xcolor}
\usepackage{amsmath, amsthm, amscd, amsfonts, amssymb}
\usepackage{graphicx}
\usepackage{stmaryrd}
\usepackage{quiver}
\usepackage{float}
\usepackage{mathtools}
\usepackage{mathrsfs}
\usepackage{multicol}
\usepackage{tikz-cd}
\usepackage[shortlabels]{enumitem}
\usepackage{stackrel}
\usepackage{cancel}
\usepackage[framemethod = TikZ]{mdframed}

%theorems
\newtheorem{theorem}{Teorema}[section]

\newtheorem{algorithm}{Algorithm}
\newtheorem{axiom}{Axioma}
\newtheorem{case}{Case}
\newtheorem{claim}{Afirmación}
\newtheorem{conclusion}{Conclusión}
\newtheorem{condition}{Condition}
\newtheorem{conjecture}[theorem]{Conjetura}
\newtheorem{corollary}[theorem]{Corolario}
\newtheorem{criterion}{Criterion}
\newtheorem{proposition}[theorem]{Proposición}
\newtheorem{lemma}[theorem]{Lema}

\theoremstyle{definition}
\newtheorem{definition}[theorem]{Definición}

\newtheorem{exercise}{Ejercicio}

\newtheorem*{notation}{Notación}
\newtheorem{problem}{Problema}[section]
\newtheorem*{obs}{Observación}
\newtheorem*{addendum}{Addendum}

\newtheorem*{tarea}{\wea{Tarea}}

\theoremstyle{remark}
\newtheorem{remark}[theorem]{Remark}
\newtheorem*{solution}{\textbf{Solución}}
\newtheorem{summary}{Summary}
\newtheorem{example}[theorem]{\textbf{Ejemplo}}

\numberwithin{equation}{section}

\renewcommand*{\proofname}{\textbf{Demostración}}

\surroundwithmdframed[
  topline=false,
  rightline=false,
  bottomline=false,
  leftmargin=\parindent,
  skipabove=\medskipamount,
  skipbelow=\medskipamount
]{proof}

\surroundwithmdframed[
  topline=false,
  rightline=false,
  bottomline=false,
  leftmargin=\parindent,
  skipabove=\medskipamount,
  skipbelow=\medskipamount
]{solution}

%commands

%Main letters
\newcommand{\A}{\mathbb{A}}
\newcommand{\B}{\mathbb{B}}
\newcommand{\C}{\mathbb{C}}
\newcommand{\D}{\mathbb{D}}
\newcommand{\F}{\mathbb{F}}
\newcommand{\N}{\mathbb{N}}
\renewcommand{\P}{\mathbb{P}}
\newcommand{\Q}{\mathbb{Q}}
\newcommand{\R}{\mathbb{R}}
\newcommand{\Z}{\mathbb{Z}}

\newcommand{\Acal}{\mathcal{A}}
\newcommand{\Bcal}{\mathcal{B}}
\newcommand{\Ccal}{\mathcal{C}}
\newcommand{\Dcal}{\mathcal{D}}
\newcommand{\Ecal}{\mathcal{E}}
\newcommand{\Fcal}{\mathcal{F}}
\newcommand{\Gcal}{\mathcal{G}}
\newcommand{\Hcal}{\mathcal{H}}
\newcommand{\Ical}{\mathcal{I}}
\newcommand{\Jcal}{\mathcal{J}}
\newcommand{\Kcal}{\mathcal{K}}
\newcommand{\Lcal}{\mathcal{L}}
\newcommand{\Mcal}{\mathcal{M}}
\newcommand{\Ncal}{\mathcal{N}}
\newcommand{\Ocal}{\mathcal{O}}
\newcommand{\Pcal}{\mathcal{P}}
\newcommand{\Qcal}{\mathcal{Q}}
\newcommand{\Rcal}{\mathcal{R}}
\newcommand{\Scal}{\mathcal{S}}
\newcommand{\Tcal}{\mathcal{T}}
\newcommand{\Ucal}{\mathcal{U}}
\newcommand{\Vcal}{\mathcal{V}}
\newcommand{\Wcal}{\mathcal{W}}
\newcommand{\Xcal}{\mathcal{X}}
\newcommand{\Ycal}{\mathcal{Y}}
\newcommand{\Zcal}{\mathcal{Z}}

%symbols
\newcommand{\vphi}{\varphi}
\newcommand{\oo}{\infty}
\newcommand{\longto}{\longrightarrow}
\newcommand{\embbed}{\hookrightarrow}
\renewcommand{\c}{\cdot}
\renewcommand{\bar}{\overline}
\newcommand{\normal}{\unlhd}
\renewcommand{\tilde}{\widetilde}
\renewcommand{\hat}{\widehat}

%special functions or sets
\newcommand{\sgn}{\mathrm{sgn}}
\newcommand{\id}{\mathrm{id}}
\newcommand{\im}{\mathrm{im}\,}
\newcommand{\op}{\mathrm{op}}
\newcommand{\GL}{\mathrm{GL}}
\newcommand{\SL}{\mathrm{SL}}
\newcommand{\Sym}{\mathrm{Sym}}
\newcommand{\Aut}{\mathrm{Aut}}
\newcommand{\Inn}{\mathrm{Inn}}

%different bracket types
\newcommand{\abs}[1]{\left| #1 \right|}
\newcommand{\norm}[1]{\left\lVert #1 \right\rVert}
\newcommand{\br}[1]{\left( #1 \right)}
\newcommand{\set}[1]{\left\{ #1 \right\}}
\newcommand{\cor}[1]{\left[ #1 \right]}
\newcommand{\gen}[1]{\left\langle #1 \right\rangle}
\newcommand{\matr}[1]{\begin{pmatrix}#1\end{pmatrix}}

\newcommand{\znz}[1]{\frac{\Z}{#1 \Z}}
\newcommand{\zmz}[1]{\Z/ #1 \Z}

%If I mess up while typing
\newcommand{\rac}[2]{\frac{#1}{#2}}
\newcommand{\olon}{\colon}
\newcommand{\ubseteq}{\subseteq}
\newcommand{\ubset}{\subset}
\newcommand{\irc}{\circ}
\renewcommand{\o}{\to}

\newcommand{\ds}{\displaystyle}

\renewcommand*\contentsname{Contenidos}

\newcommand{\wea}[1]{\textcolor{BlueViolet}{\textbf{\emph{#1}}}}
\newcommand{\fecha}[1]{\begin{flushright} \underline{#1} \end{flushright}}
\newcommand{\dem}[1]{\textbf{Demostración (#1)}}

%title
\title{MPG3100 -- Álgebra I}
\author{Felipe Inostroza \\ Prof. Sebastián Barbieri}
\date{Primer Semestre 2026}

\begin{document}

\maketitle
\tableofcontents

\newpage
\fecha{Clase 1. Lunes 9 de Marzo}
\section{Grupos}

\begin{definition}
  Un \wea{grupo} es un par $(G,\mathrm{op})$ donde $G$ es un conjunto y $\mathrm{op} \colon G \times G \to G$ tal que si denotamos $g\cdot h = \mathrm{op}(g,h)$, $\forall g,h \in G$, entonces
  \begin{enumerate}[(1)]
    \item $\forall x,y,z \in G$, $(x \cdot y) \cdot z = x \cdot (y \cdot z)$;
    \item $\exists e \in G$ tal que $\forall x \in G$, $e \cdot x = x \cdot e = x$.
    \item $\forall x \in G$, $\exists y \in G$ tal que $x \cdot y = y \cdot x = e$.
  \end{enumerate}
\end{definition}

\begin{definition}
  $(G,\op)$ se dice \wea{abeliano} si $\op(x,y) = \op(y,x)$, $\forall x,y \in G$.
\end{definition}

\begin{obs}
  En todo grupo, el neutro es único.
\end{obs}

\begin{notation}
  El neutro de un grupo $(G,\cdot)$ se denota:
  \begin{itemize}
    \item $e, e_G$,
    \item $1, 1_G$,
    \item $0, 0_G$ (exclusivamente si $G$ es abeliano)
    \item $\id$.
  \end{itemize}
  La operación se suele denotar por
  \begin{itemize}
    \item $\op(x,y) = x \cdot y$ (en general para grupos no abelianos)
    \item $\op(x,y) = x+y$ (exclusivamente para grupos abelianos)
  \end{itemize}
  Si $xy = yx = e$, escribimos $y = x^{-1}$.
  Si $x+y=y+x=e$, escribimos $y = -x$.
\end{notation}

\begin{definition}
  Sea $G$ grupo. Un subconjunto $H \subset G$ se dice \wea{subgrupo} (denotado $H \leqslant G$) si $H$ con la operación heredada de $G$ es grupo.
\end{definition}

\begin{proposition}
  Sea $H \subset G$ no vacío. Entonces $H \leqslant G \iff \forall h,h' \in H, h'h^{-1} \in H$.
\end{proposition}

\begin{example}
  $(\Z,+) \leqslant (\Q,+) \leqslant (\R,+) \leqslant (\C,\leqslant)$.
\end{example}

\begin{example}
  Si $(V,+,\cdot)$ es un espacio vectorial, entonces $(V,+)$ es un grupo abeliano. Por ejemplo $(\R^n,+)$ y $(\R[x],+)$ son grupos abelianos.
\end{example}

\begin{example}
  Si $M_n(\R)$ es el conjunto de matrices de $n \times n$ sobre $\R$, entonces
  \[\GL_n(\R) = \{A \in M_n(\R) \colon \det(A) \neq 0\}\]
  es un grupo no abeliano.
\end{example}

\begin{example}
  $\SL_n(\R) = \{A \in \GL_n(\R) \colon \det(A) = 1\}$ es un subgrupo de $\GL_n(\R)$.
\end{example}

\begin{example}
  Dado un conjunto $X$, consideramos el \wea{grupo simétrico}
  \[\Sym(X) = \{f \colon X \to X \text{ biyectivas}\}.\]
  Tenemos que $(\Sym(X),\circ)$ es un grupo.
\end{example}

\begin{example}
  $S^1 = \{z \in \C \colon |z| = 1\}$ con la multiplicación.

  Fijemos $n \geq 1$ entero.
  \[\Z_n = \{z \in S^1 \colon z^n = 1\}.\]
  Se tiene que $\Z_n \leqslant S^1$.
\end{example}

\begin{definition}
  Sean $G,H$ grupos. Un \wea{homomorfismo} es una función $$\phi \colon G \to H$$ tal que $\forall x, y \in G$, $$\phi(x\c y) = \phi(x) \c \phi(y).$$
\end{definition}

\begin{obs}
  $\phi(1_G) = 1_H$, $\phi(g^{-1}) = (\phi(g))^{-1}$.
\end{obs}

\begin{definition}
  Sea $\phi \colon G \to H$ homomorfismo de grupos. La \wea{imagen} y el \wea{kernel} de $\phi$ son (resp.)
  \begin{align*}
    \im(\phi) &= \{h \in H \colon \exists g \in G, \phi(g) = h\}, \\
    \ker(\phi) &= \{g \in G \colon \phi(g) = 1_H\}.
  \end{align*}
\end{definition}

\begin{proposition}
  $\im(\phi) \leqslant H$, $\ker(\phi) \leqslant G$.
\end{proposition}
\begin{proof}
  \
  \begin{itemize}
    \item \boxed{\im\phi} Primero notemos que $\phi(e_G) = e_H \in \im(\phi)$, por lo que $\im(\phi) \neq \varnothing$. Ahora, sean $h_1,h_2 \in \im(\phi)$, entonces existen $g_1,g_2 \in G$ tales que $h_i = \phi(g_i)$. Luego $$h_1 h_2^{-1} = \phi(g_1) (\phi(g_2))^{-1} = \phi(g_1) \phi(g_2^{-1}) = \phi(g_1 g_2^{-1}) \in \im(\phi)$$ y como $h_1,h_2$ son arbitrarios, tenemos que $\im(\phi) \leqslant H$.
    \item \boxed{\ker(\phi)} Sabemos que $\phi(e_G) = e_H$, por lo que $e_G \in \ker(\phi)$, y entonces $\ker(\phi) \neq \varnothing$. Sean $g,h \in \ker(\phi)$, entonces $\phi(g) = \phi(h) = e_H$. Luego $$\phi(gh^{-1}) = \phi(g) (\phi(h))^{-1} = e_H \c e_H = e_H$$ y entonces $gh^{-1} \in \ker(\phi)$. Por lo tanto $\ker(\phi) \leqslant G$.
  \end{itemize}
\end{proof}

%% demostrar
\begin{exercise}
  Un homomorfismo $\phi \colon G \to H$ es inyectivo si y solo si $\abs{\ker(\phi)} = 1$.
\end{exercise}
\begin{proof}
  \boxed{\implies}
  Supongamos que $\phi \colon G \to H$ es un homomorfismo inyectivo. Sea $x \in \ker(\phi)$, entonces $\phi(x) = e_G = \phi(e_G)$, y como $\phi$ es inyectivo, tenemos que $x = e_G$. Por lo tanto $\ker(\phi) = \{e_G\}$, y entonces $\abs{\ker(\phi)} = 1$.

  \boxed{\impliedby}
  Supongamos que $\abs{\ker(\phi)} = 1$. Sean $g,h \in G$ tales que $\phi(g) = \phi(h)$. Entonces
  $$e_H = \phi(g) \c (\phi(h))^{-1} = \phi(gh^{-1})$$
  por lo que $gh^{-1} \in \ker(\phi)$. Como $e_G \in \ker(\phi)$, y $|\ker(\phi)| = 1$, tenemos que $\ker(\phi) = \{e_G\} \ni gh^{-1}$. Luego $gh^{-1} = e_G \implies g=h$. Por lo tanto $\phi$ es inyectivo.
\end{proof}

\begin{definition}
  El \wea{grupo trivial} es un grupo con un único elemento. Se denota por $G = 1$.
\end{definition}

\begin{example}
  Sea $n \geq 1$ entero, y consideremos
  \begin{align*}
    \psi_n \colon S^1 &\to S^1 \\
    z &\mapsto z^n
  \end{align*}
  Esta función es homomorfismo, con $\im(\psi_n) = S^1$, $\ker(\psi_n) = \Z_n$.
\end{example}

\begin{example}
  Toda transformación lineal $T \colon \R^n \to \R^n$ es homomorfismo.
  \begin{align*}
    \phi \colon \Z &\to \GL_2(\Z) \\
    n &\mapsto \matr{1&n\\0&1}
  \end{align*}
\end{example}

\begin{notation}
  Sea $\phi \colon G \o H$ homomorfismo. Diremos que
  \begin{itemize}
    \item $\phi$ es \wea{monomorfismo} si $\phi$ es inyectivo.
    \item $\phi$ es \wea{epimorfismo} si $\phi$ es sobreyectivo.
    \item $\phi$ es \wea{isomorfismo} si $\phi$ es biyectivo.
    \item $\phi$ es \wea{endomorfismo} si $G=H$.
    \item $\phi$ es \wea{automorfismo} si $\phi$ es biyectiva y $G=H$.
  \end{itemize}
  Si existe un isomorfismo entre $G$ y $H$, decimos que son \wea{isomorfos}, denotado por $G \simeq H$.
\end{notation}

%% demostrar
\begin{example}
  Sea $K = \set{\matr{x&-y\\y&x} \in M_2(\R)}$ con la suma. Tenemos que $(K,+) \simeq (\C,+)$.

  Si $K^\ast = K \setminus \set{\matr{0&0\\0&0}}$, entonces $(K^\ast,\c) \simeq (\C \setminus \{0\}, \c)$.
\end{example}

\begin{example}
  $(\zmz n, +) \simeq (\Z_n, \cdot)$.
\end{example}

\begin{exercise}
  Entre los siguientes grupos, determine cuales son isomorfos entre si y cuales no.
  \[(\Z,+), \quad (\Q,+), \quad (\R,+), \quad (\C,+).\]
\end{exercise}
\begin{proof}\
  \begin{itemize}
    \item \boxed{\Z \not\simeq \Q} Supongamos que existe un isomorfismo $\phi \colon \Z \to \Q$. Sea $q = \phi(1)$, luego existe $u \in \Z$ tal que $\phi(u) = q/2$. Entonces $$\phi(u+u) = \phi(u) + \phi(u) = q$$ y entonces $u+u = 1$. $\to\gets$
    \item \boxed{\Z \not\simeq \R} Recordemos que $|\Z| \neq |\R|$, es decir, no existe biyección entre $\Z$ y $\R$, y por lo tanto no puede existir un isomorfismo $\Z \to \R$.
    \item \boxed{\Z \not\simeq \C} Mismo argumento.
    \item \boxed{\Q \not\simeq \R} Mismo argumento.
    \item \boxed{\Q \not\simeq \C} Mismo argumento.
    \item \boxed{\R \simeq \C} Notemos que $\R$ y $\C$ son espacios vectoriales sobre $\Q$, y ambos tienen dimensión $\aleph_0$ (con base de Hamel), por lo que existe un isomorfismo de espacios vectoriales $\phi \colon \R \to \C$. La restricción de este isomorfismo con la suma nos da un isomorfismo de grupos.
  \end{itemize}
\end{proof}

\begin{definition}
  Sea $G$ un grupo y $g \in G$. La \wea{conjugación} es la aplicación
  \begin{align*}
    C_g \colon G &\to G \\
    h &\mapsto ghg^{-1}.
  \end{align*}
  $C_g$ es un automorfismo.
\end{definition}

\begin{notation}
  Dado un grupo $G$,
  $$\Aut(G) = \{\varphi \colon G \to G \text{ automorfismo}\}$$
  con la composición es el \wea{grupo de automorfismos de $G$}.
  $$\Inn(G) = \{\varphi \colon G \to G \mid \exists g \in G, \varphi = C_g\}$$
  con la composición es el \wea{grupo de automorfismos internos}.
\end{notation}

\begin{obs}
  $\Inn(G) \leqslant \Aut(G)$.
\end{obs}

\begin{example}
  Si $(G,+)$ es abeliano, entonces $\Inn(G) \simeq 1$.
\end{example}

\begin{exercise}
  Pruebe que $$\Aut(\Z^d,+) \simeq (\GL_d(\Z), \c).$$
\end{exercise}

%% demostrar cosas

\begin{exercise}
  Mostrar que $(\Q,+) \not\cong (\Q^2,+)$ y $(\Z,+) \not\cong (\Z^2,+)$
\end{exercise}

\newpage
\fecha{Clase 2. Miércoles 11 de Marzo}

\begin{proposition}
  Dados $X$ e $Y$ conjuntos, entonces
  \[\Sym(X) \cong \Sym(Y) \iff |X| = |Y|\]
\end{proposition}
\begin{proof}
  \
  \begin{itemize}
    \item \boxed{\impliedby} Supongamos que $|X| = |Y|$. Entonces existe $f \colon X \to Y$ biyectiva. Definimos $$\phi \colon \Sym(X) \to \Sym(Y)$$ dada por $\phi(\sigma) = f \circ \sigma \circ f^{-1}$ para todo $\sigma \in \Sym(X)$.
  \end{itemize}
\end{proof}
%% demostrar la otra implicancia

\begin{notation}
  \begin{align*}
    S_n &= \Sym(\{1,\dots,n\}) \\
    S_\oo &= \Sym(\Z)
  \end{align*}
  $\Sym_0(\Z)$ es el grupo de permutaciones de $\Z$ de soporte finito.
\end{notation}

\begin{obs}
  $\abs{S_n} = n!$.
\end{obs}

\begin{definition}
  Decimos que $\sigma \in S_n$ es un \wea{ciclo de largo $k$} si existen elementos $a_1,\dots,a_k \in \{1,\dots,n\}$ tales que
  \begin{itemize}
    \item $\sigma(a_i) = a_{i+1}, \quad \forall i \in \{1,\dots,k-1\}$;
    \item $\sigma(a_k) = a_1$;
    \item $\sigma(x) = x, \quad \forall x \in \{1,\dots,n\} \setminus \{a_1,\dots,a_k\}$.
  \end{itemize}
\end{definition}

\begin{notation}
  El ciclo anterior se denota
  \[\sigma = (a_1,a_2,a_3,\dots,a_k).\]
\end{notation}

\begin{obs}
  La notación no es única.
\end{obs}

\begin{example}
  En $S_5$, tenemos que
  \[(1,2)\c(2,3,5) = (5,1,2,3).\]
\end{example}

\begin{obs}
  Si $\tau = (a_1,\dots,a_k)$ y $\sigma = (b_1,\dots,b_l)$ son tales que
  \[\{a_1,\dots,a_k\} \cap \{b_1,\dots,b_l\} = \varnothing.\]
  Entonces $\tau \circ \sigma = \sigma \circ \tau$.
\end{obs}

\begin{obs}
  Toda permutación en $S_n$ se escribe como producto finito de ciclos disjuntos.
\end{obs}

\begin{definition}
  Una \wea{transposición} es un ciclo de largo 2.
\end{definition}

\begin{obs}
  Todo ciclo en $S_n$ se escribe como producto de transposiciones:
  \[(a_1,\dots,a_k) = (a_1,a_k)(a_1,a_{k-1}) \dots (a_1,a_3)(a_1,a_2).\]
\end{obs}

\begin{obs}
  Toda permutación en $S_n$ es producto de transpocisiones.
\end{obs}

\begin{definition}
  Sea $G$ grupo y $S \subset G$. El \wea{subgrupo de $G$ generado por $S$} es
  \[\gen{S} = \{g \in G \colon \exists n \geq 0 \colon g = s_1,\dots,s_n, \quad s_1,\dots,s_n \in S \cup S^{-1}\}.\]
  Diremos que $S \subset G$ es \wea{generador} si $\gen{S} = G$.
  Diremos que $G$ es \wea{finitamente generado} si $\exists S \subset G$ finito y generador.
\end{definition}

\begin{obs}
  Si $G$ es finitamente generado, entonces $G$ es contable (finito o numerable).
\end{obs}

\begin{obs}
  $S_n$ es generado por las transposiciones.
\end{obs}

\begin{exercise}
  Mostrar que $\Sym_0(\Z)$ no es finitamente generado.
\end{exercise}

\begin{theorem}
  Si $\sigma \in S_n$ se puede escribir como producto de una cantidad par de transposiciones, entonces no es posible escribirla como producto de una cantidad impar de transposiciones.
\end{theorem}
\begin{proof}
  Primero demostrar para la identidad. %% demostrar
  Si
  \[\tau = t_1 t_2 \dots t_k = r_1 r_2 \dots r_l\]
  con $t_i, r_i$ transposiciones, entonces
  \[\id = t_1 \dots t_k r_l^{-1} \dots r_1^{-1}
  \implies k \equiv l \pmod 2.\]
\end{proof}

\begin{definition}
  Consideremos la función $\phi \colon S_n \to \Z$ dada por
  \begin{itemize}
    \item $\phi(\sigma) = 1$ si $\sigma$ es permutación par;
    \item $\phi(\sigma) = -1$ si $\sigma$ es permutación impar.
  \end{itemize}
  El \wea{grupo alternante} es
  \[A_n = \ker(\phi) = \{\sigma \in S_n \colon \sigma \text{ es par}\}.\]
\end{definition}

\begin{definition}
  El \wea{grupo diedral} $D_n$ es el grupo de simetrias el polígono regular de $n$ lados.
\end{definition}

\begin{example}
  $D_3$ es el grupo de simetrías de un triángulo equilátero. Si denotamos los vértices como 1, 2 y 3, entonces $D_3$ permuta estos números entre si, por lo que $D_3 \leqslant S_3$, y
  \[D_3 = \{\id, (123), (132), (13), (12), (23)\}.\]
  Similarmente, $D_4 \leqslant S_4$, pero no escribiré sus elementos.
\end{example}

\newpage
\subsection{Clases laterales}

\begin{definition}
  Sea $G$ un grupo, $H \leqslant G$. Una \wea{clase lateral} (\wea{izquierda} o \wea{derecha}) es
  \[gH = \{gh \colon h \in H\}, \quad Hg = \{hg \colon h \in H\}\]
  para algún $g \in G$.
\end{definition}

\begin{proposition}
  Las clases laterales de $H \leqslant G$ sorman una partición de $H$:
  \begin{enumerate}[(i)]
    \item $\bigcup_{g\in G} gH = G$;
    \item $gH \cap g'H \neq \varnothing \implies gH = g'H$.
  \end{enumerate}
\end{proposition}
\begin{proof}
  $\bigcup_{g \in G} gH = G$ es trivial.

  Supongamos que $gH \cap g' H \neq \varnothing$, entonces $\exists h,h' \in H$ tales que $gh = g'h'$. Luego $g' = gh(h')^{-1} \in gH$. Luego $\forall h'' \in H$, tenemos que $g'h'' \in gH$, y entonces $g'H \subseteq gH$. Análogo para $gH \subseteq g'H$.
\end{proof}

\begin{obs}
  Supongamos que $H \leqslant G$ y $\abs{H} < \oo$, entonces $\forall g \in G$, $|gH| = |H|$. Si $G$ es finito, entonces
  \[G = \bigcup_{g \in L} gH \implies \abs{G} = \abs{L} \c \abs{H}.\]%% agregar punto dentro del bigcup
\end{obs}

\begin{definition}
  Decimos que $L \subset G$ es un \wea{conjunto de representantes izquierdos} (\wea{derechos}) de $H \leqslant G$ si
  \[G = \bigcup_{g \in L} gH = \bigcup_{g \in L} Hg.\]%% agregar punto dentro de los bigcups
\end{definition}

\begin{theorem}
  [Lagrange]
  Sea $G$ un grupo finito. Si $H \leqslant G$, entonces $|H|$ divide a $|G|$.
\end{theorem}

\begin{definition}
  Un subgrupo $G \leqslant G$ se dice \wea{normal} y se denota $H \lhd G$ si $\forall h \in H$ y $g \in G$,
  \[ghg^{-1} \in H.\]
\end{definition}

\begin{proposition}
  $H \lhd G \iff \forall g \in G, \ gH = Hg$.
\end{proposition}

\begin{obs}
  Si $G$ es abeliano, entonces $H \leqslant G \implies H \lhd G$.
\end{obs}

\begin{example}
  Todos los subgrupos de $(\Z,+)$ son de la forma $n\Z$ con $n \geq 1$. Como $\Z$ es abeliano, estos son subgrupos normales. Las clases laterales son de la forma $k + n\Z$ con $k \in \{0,\dots,n-1\}$. Notemos que
  \[(k + n\Z) + (k'+n\Z) = (k+k' \pmod n + n\Z)\]
\end{example}



%% cambiar todos los lhd por parecido a leqslant


%\fecha{Clase 3.}
%\fecha{Clase 4.}
%\fecha{Clase 5.}
%\fecha{Clase 6.}
%\fecha{Clase 7.}
%\fecha{Clase 8.}
%\fecha{Clase 9.}
%\fecha{Clase 10.}
%\fecha{Clase 11.}
%\fecha{Clase 12.}
%\fecha{Clase 13.}
%\fecha{Clase 14.}
%\fecha{Clase 15.}
%\fecha{Clase 16.}
%\fecha{Clase 17.}
%\fecha{Clase 18.}
%\fecha{Clase 19.}
%\fecha{Clase 20.}
%\fecha{Clase 21.}
%\fecha{Clase 22.}
%\fecha{Clase 23.}
%\fecha{Clase 24.}
%\fecha{Clase 25.}
%\fecha{Clase 26.}
%\fecha{Clase 27.}
%\fecha{Clase 28.}
%\fecha{Clase 29.}
%\fecha{Clase 30.}


\end{document}